\documentclass[10pt,a4paper]{amsart}
\usepackage[margin=19mm]{geometry}
\usepackage{amsmath,amssymb,mathtools}
\usepackage[scr=rsfs,cal=euler]{mathalfa}
\usepackage{xfrac}
\usepackage[unicode,hidelinks,bookmarks=false]{hyperref}
\newcommand{\ie}{i.e.\ }
\newcommand{\cf}{cf.\ }
\newcommand{\resp}{respectively\ }
\newcommand{\Subsection}{\subsection}
\providecommand{\nobreakdash}{\nobreak\hbox{-}}
\setlength{\emergencystretch}{2em}
\multlinegap0pt
\usepackage{amssymb}
\usepackage{bm}
\usepackage{xcolor}
\newtheorem{proposition}{Proposition}
\title{Transient Acceleration and Cross-Dissipation Interference in Fisher-Regularized Wasserstein Gradient Flows}
\author{Michael Farmer, Abhinav Kochar,, Yugyung Lee}
\date{Source: arXiv:2603.12285v2 (CC BY 4.0)}
\begin{document}
\maketitle

\noindent\emph{Source and licence.} Transient Acceleration and Cross-Dissipation Interference in Fisher-Regularized Wasserstein Gradient Flows. Version arXiv:2603.12285v2 (CC BY 4.0), distributed by arXiv under the Creative Commons Attribution 4.0 International licence, http://creativecommons.org/licenses/by/4.0/, as recorded in the arXiv licence metadata for this exact version and shown on the abstract page https://arxiv.org/abs/2603.12285v2. Full licence notice: \texttt{licenses/arxiv-2603.12285-CC-BY-4.0.txt}.\par\smallskip

\noindent\textit{Editorial scope.} Counted unit: the article's proposition proposition (source0.tex lines 479--506). The article's complete original proof of it is delivered with it (source0.tex lines 508--547, one \texttt{proof} environment of 40 lines). No mathematics is written, repaired or re-derived: the statement and the proof are the article's own lines, in the article's own notation, and no retained line is substituted or re-flowed.  No further source-local context is needed: every cross-reference of every retained line resolves inside the two delivered spans.  Nothing was omitted from any retained span, so the package carries no omission record.  No retained line contains a citation, so the wrapper carries no bibliography and no bibliography entry.  No retained line contains a figure environment, an image-inclusion command or drawing code, so no image is delivered and the package declares no asset dependency.  Editorial formatting only, in addition: an AMS \texttt{amsart} preamble that restates the article's own declarations and macros used by the retained lines, namely the environments \texttt{proposition} on separate counters, together with the standard packages those lines need.  The retained environments print in this document as Proposition 1 in order of appearance, whereas the article prints the same items with the numbering of its own full document.  The complete native source of the article as distributed in the arXiv e-print for this exact version is bundled unchanged as \texttt{sources/196268/source0.tex}.
\begin{proposition}[Bounded acceleration-window crossing time]
For Gaussian initial data with variance
\[
u_0<1,
\]
the transition time required to reach the critical scale
\[
u=1
\]
is given by
\[
T_{\mathrm{acc}}(u_0)
=
\int_{u_0}^{1}
\frac{u\,du}
{2u(1-u)+\varepsilon}.
\]
Moreover, for fixed $\varepsilon>0$, the transition time is finite and monotonically decreasing as the initial variance $u_0$ increases. In particular,
\[
T_{\mathrm{acc}}(u_0)
\le
\int_{0}^{1}
\frac{u\,du}
{2u(1-u)+\varepsilon}
=: C_\varepsilon < \infty.
\]
Hence, the duration of the cooperative acceleration regime is uniformly bounded by a constant depending only on the Fisher regularization strength $\varepsilon$.
\end{proposition}

\begin{proof}
The cooperative acceleration regime terminates when the trajectory reaches the critical scale $u=1$. From the reduced Gaussian-manifold dynamics,
\[
\dot u
=
2(1-u)
+
\frac{\varepsilon}{u},
\]
we obtain
\[
dt
=
\frac{u\,du}
{2u(1-u)+\varepsilon}.
\]
Integrating from $u_0$ to $1$ gives
\[
T_{\mathrm{acc}}(u_0)
=
\int_{u_0}^{1}
\frac{u\,du}
{2u(1-u)+\varepsilon}.
\]
For $\varepsilon>0$, the denominator satisfies
\[
2u(1-u)+\varepsilon>0
\]
on $0<u\leq 1$, so the integral is finite. Furthermore,
\[
T_{\mathrm{acc}}(u_0)
\le
\int_{0}^{1}
\frac{u\,du}
{2u(1-u)+\varepsilon}
=
C_\varepsilon,
\]
which establishes a uniform upper bound independent of the initial variance. Monotonicity with respect to $u_0$ follows directly from the integral representation.
\end{proof}

\end{document}
